% Part: first-order-logic
% Chapter: introduction
% Section: formulas

\documentclass[../../../include/open-logic-section]{subfiles}

\begin{document}

\olfileid[hi]{fol}{int}{fml}

\section{\usetoken{P}{formula}}

प्रथम-क्रम तर्क के !!{sentence}s को कठोरता से निर्दिष्ट करने और !!{variable}s
के उपयोग से उठने वाली समस्याओं को सँभालने के लिए हम निम्न पद्धति अपनाएँगे।
पहले व्यंजकों के एक \emph{अलग} समुच्चय, अर्थात् !!{formula}s के समुच्चय, को परिभाषित
करेंगे। इसके बाद हम उनमें !!{variable}s की भूमिका पर विचार कर सकेंगे और कुछ
अन्य धारणाओं---अर्थात् ``मुक्त'' और ``बद्ध'' !!{variable}s---के आधार पर
परिभाषित कर सकेंगे कि !!a{sentence} क्या है (अर्थात् ऐसा !!a{formula} जिसमें
कोई मुक्त !!{variable} न हो)। हम ऐसा केवल इसलिए नहीं करते कि इससे
``!!{sentence}'' की परिभाषा अधिक सुगम हो जाती है, बल्कि इसलिए भी कि संतुष्टि
की अर्थवैज्ञानिक धारणा को परिभाषित करने की हमारी विधि में यह निर्णायक होगा।

अब किसी सरल प्रथम-क्रम भाषा के लिए ``!!{formula}'' परिभाषित करें। इस भाषा
में केवल एक !!{predicate}~$\Obj P$, एक !!{constant}~$\Obj a$ और तार्किक
चिह्न $\lnot$, $\land$ तथा~$\lexists$ हैं। हमारी पूरी परिभाषाएँ कहीं अधिक
सामान्य होंगी: उनमें अनंततः अनेक !!{predicate}s और !!{constant}s की अनुमति
होगी। वास्तव में हम ऐसे !!{function}s पर भी विचार करेंगे जिन्हें
!!{constant}s और !!{variable}s के साथ जोड़कर ``पद'' बनाए जा सकते हैं। अभी
$\Obj a$ और चर ही हमारे पद होंगे। हमें अनंततः अनेक !!{variable}s की आवश्यकता है।
औपचारिक रूप से हम $\Obj v_0$, $\Obj
v_1$, \dots, चिह्नों को चर मानेंगे।

\begin{defn}
\emph{!!{formula}s} के समुच्चय~$\Frm$ को निम्न प्रकार परिभाषित किया जाता है:
\begin{enumerate}
\item\ollabel{fmls-atom} $\Atom{\Obj P}{\Obj a}$ और $\Atom{\Obj
  P}{\Obj v_i}$ को !!{formula}s में गिना जाता है ($i \in \Nat$)।

\tagitem{prvNot}{\ollabel{fmls-not}यदि $!A$ कोई !!a{formula} है, तो $\lnot !A$
  भी !!a{formula} है।}{}

\tagitem{prvAnd}{यदि $!A$ और $!B$ को !!{formula}s में गिना जाता है, तो $(!A \land
  !B)$ कोई !!a{formula} है।}{}

\tagitem{prvEx}{\ollabel{fmls-ex}यदि $!A$ कोई !!a{formula} और $x$ कोई
  !!a{variable} है, तो $\lexists[x][!A]$ कोई !!a{formula} है।}{}

\tagitem{limitClause}{\ollabel{fmls-limit}इसके अतिरिक्त कुछ भी !!a{formula} नहीं है।}{}
\end{enumerate}
\end{defn}

\olref{fmls-atom} बताती है कि $\Atom{\Obj P}{\Obj a}$ और
$\Atom{\Obj P}{\Obj v_i}$ को, किसी भी $i \in
\Nat$ के लिए, !!{formula}s में गिना जाता है। इन्हें तथाकथित
\emph{परमाण्विक} !!{formula}s के वर्ग में रखा जाता है और ये हमें आरंभ-बिंदु
देते हैं। अन्य खंड पहले से बने सूत्रों से नए !!{formula}s के निर्माण की विधियाँ देते हैं। उदाहरणार्थ,
\olref{fmls-not} से $\lnot \Atom{\Obj P}{\Obj v_2}$ कोई !!a{formula} है,
क्योंकि \olref{fmls-atom} के अनुसार $\Atom{\Obj P}{\Obj v_2}$ पहले ही
!!a{formula} है। फिर \olref{fmls-ex} से
$\lexists[\Obj v_2][\lnot \Atom{\Obj P}{\Obj v_2}]$ भी !!a{formula} है,
और यह क्रम आगे चलता है।
\olref{fmls-limit} बताती है कि \emph{केवल} इसी ढंग से बनाई जा सकने वाली
शृंखलाओं को !!{formula}s की श्रेणी में रखा जाता है। विशेषतः,
$\lexists[\Obj v_0][\Atom{\Obj P}{\Obj a}]$ और $\lexists[\Obj
  v_0][\lexists[\Obj v_0][\Atom{\Obj P}{\Obj a}]]$ को !!{formula}s की श्रेणी में
\emph{रखा जाता है}, जबकि अनावश्यक बाहरी कोष्ठकों के
कारण $(\lnot \Atom{\Obj P}{\Obj a})$ को नहीं।

!!{formula}s को इस प्रकार परिभाषित करना \emph{आगमनात्मक परिभाषा} कहलाता है।
इससे हम आगमन द्वारा प्रमाण के एक रूप, \emph{संरचनात्मक आगमन}, का उपयोग करके
!!{formula}s के बारे में बातें सिद्ध कर सकते हैं। इनकी सामान्य चर्चा
\olref[mth][ind][idf]{sec} और \olref[mth][ind][sti]{sec} में है; आगे के
प्रमाणों में जाने से पहले आपको उन्हें दोहरा लेना चाहिए। मूल विचार यह है कि
यदि आप कोई बात सभी !!{formula}s के लिए सिद्ध करना चाहते हैं, तो पहले दिखाते
हैं कि वह परमाण्विक !!{formula}s के लिए सत्य है; फिर दिखाते हैं कि वह किसी
!!{formula}~$!A$ (और~$!B$) के लिए सत्य \emph{हो}, तो $\lnot !A$,
$(!A \land !B)$ और $\lexists[x][!A]$ के लिए \emph{भी} सत्य है। उदाहरणार्थ,
इससे वह बात
$\lexists[\Obj v_2][\lnot \Atom{\Obj P}{\Obj v_2}]$ के लिए सिद्ध होती है:
पहले भाग से वह परमाण्विक !!{formula}~$\Atom{\Obj
P}{\Obj v_2}$ के लिए सत्य है। फिर दूसरे भाग से वह $\lnot \Atom{\Obj
P}{\Obj v_2}$ के लिए और पुनः
$\lexists[\Obj v_2][\lnot \Atom{\Obj P}{\Obj v_2}]$ के लिए सत्य है। चूँकि
!!{formula}s का पूरा वर्ग परमाण्विक !!{formula}s से आगमनात्मक रूप से उत्पन्न होता है,
यह विधि उनमें से प्रत्येक के लिए काम करती है।

\end{document}
